\SourceLeaf{318}%
für die Kugeloberfläche angegebenen entsprechen und in diese übergehen, wenn man die Exzentrizität \(=0\) setzt oder die beiden Halbachsen der erzeugenden Ellipse einander gleich und \(=1\) setzt.

\Needspace{17\baselineskip}
Wie oben in Artikel 3 bezeichne ich die Breite des Normalparallelkreises, der der konformen Abbildung auf die Kugeloberfläche zugrunde liegt, für die Ellipsoidoberfläche mit \(P\) und für die Kugeloberfläche mit \(Q\). Zugleich nehme ich an, dass dieser Normalparallelkreis so gewählt ist, dass \(Q\) dem arithmetischen Mittel der Breiten der beiden betreffenden Punkte auf der Kugeloberfläche entspricht. Die Breiten selbst seien \(Q+\frac12q\) und \(Q-\frac12q\). Ferner sollen \(a,A,\alpha,e,\varphi,\theta\) ihre Bedeutungen aus den Artikeln 2, 3, 4 ff. der ersten Abhandlung behalten. Dabei bezeichnet
\begin{list}{}{\setlength{\leftmargin}{1em}\setlength{\itemindent}{0pt}\setlength{\itemsep}{0pt}\setlength{\parsep}{0pt}}
\item \(a\) die große Halbachse des Ellipsoids oder den Radius des Äquators,
\item \(A\) den Radius der Kugel,
\item \(1:\alpha\) das konstante Verhältnis der Längenunterschiede auf dem Ellipsoid zu den entsprechenden Längenunterschieden auf der Kugel,
\item \(e=\sin\varphi\) die Exzentrizität der erzeugenden Ellipse,
\item \(\sin\theta=e\sin P\).
\end{list}

\indent Den Großkreisbogen zwischen den beiden Punkten auf der Kugeloberfläche bezeichne ich mit \(s\), die Azimute dieses Bogens am ersten und zweiten Endpunkt mit \(U+\frac12u\) und \(U-\frac12u\pm180^0\). Berücksichtigt man außerdem, dass der Längenunterschied zwischen beiden Punkten \(=\alpha l\) ist, so erhält man zunächst die vier exakten Formeln
\[
\begin{aligned}
\sin\tfrac12s\mathord{.}\cos U&=\cos\tfrac12\alpha l\mathord{.}\sin\tfrac12q\\
\sin\tfrac12s\mathord{.}\sin U&=\sin\tfrac12\alpha l\mathord{.}\cos Q\\
\cos\tfrac12s\mathord{.}\cos\tfrac12u&=\cos\tfrac12\alpha l\mathord{.}\cos\tfrac12q\\
\cos\tfrac12s\mathord{.}\sin\tfrac12u&=\sin\tfrac12\alpha l\mathord{.}\sin Q
\end{aligned}
\]
und daraus die Näherungsformeln
\[
\begin{aligned}
q&=s\mathord{.}\cos U(1+\tfrac1{24}qq-\tfrac1{24}ss+\tfrac18\alpha\alpha ll)\quad .\ .\ .\ .\ .\ .\ .\quad(1)\\
\alpha l&=s\mathord{.}\frac{\sin U}{\cos Q}(1-\tfrac1{24}ss+\tfrac1{24}\alpha\alpha ll)\quad .\ .\ .\ .\ .\ .\ .\ .\ .\ .\quad(2)\\
u&=s\mathord{.}\sin U\mathord{.}\operatorname{tang}Q(1+\tfrac1{12}ss+\tfrac1{24}uu)\quad .\ .\ .\ .\ .\ .\ .\ .\quad(3)
\end{aligned}
\]
Es braucht kaum erwähnt zu werden, dass \(l,q,s,u\) in diesen drei Gleichungen in Einheiten des Radius ausgedrückt sind. Man sieht leicht, dass die Gleichungen
